\documentclass[12pt]{article}
\usepackage[utf8]{inputenc}
\usepackage[spanish]{babel}
\usepackage{geometry}
\usepackage{enumitem}
\usepackage{tabularx}
\usepackage{booktabs}
\usepackage{graphicx}
\usepackage{listings}
\usepackage{xcolor}
\usepackage{tikz}
\usetikzlibrary{shapes,arrows,positioning,calc}

\geometry{a4paper, margin=1.5cm}

% Configuración de listados para código MIPS
\lstset{
    language=[mips]Assembler,
    basicstyle=\ttfamily\small,
    keywordstyle=\color{blue},
    commentstyle=\color{green!50!black},
    stringstyle=\color{red},
    numbers=left,
    numberstyle=\tiny\color{gray},
    stepnumber=1,
    numbersep=5pt,
    backgroundcolor=\color{white},
    frame=single,
    frameround=tttt,
    rulecolor=\color{black},
    breaklines=true,
    breakatwhitespace=true,
    tabsize=4,
    showstringspaces=false
}

\title{\textbf{Informe de Laboratorio 4 \\ Interrupciones y Rutinas de Servicio en MIPS32}}
\author{
    Samuel Cordero \\ C.I V-31.678.592 \\
    Dervis Martínez \\ C.I V-31.456 \\
    \vspace{1cm}
    Universidad de Carabobo \\
    Facultad de Ciencias y Tecnología (FACYT) \\
    Departamento de Computación \\
    Arquitectura del Computador \\
    Julio 2025
}
\date{}

\begin{document}

\maketitle

\section*{1. Ciclo de Interrupción de Hardware}
El ciclo de interrupción de hardware sigue una secuencia específica desde que ocurre el evento hasta que se atiende:

\begin{enumerate}
    \item \textbf{Ocurrencia del evento}: Un dispositivo periférico (teclado, temporizador, etc.) detecta un evento que requiere atención.
    \item \textbf{Solicitud de interrupción (IRQ)}: El dispositivo activa la línea de interrupción correspondiente.
    \item \textbf{Detección por la CPU}: Al final del ciclo de instrucción, la CPU verifica si hay interrupciones pendientes.
    \item \textbf{Habilitación de interrupciones}: Si las interrupciones están habilitadas (bit IE en registro Status):
    \begin{itemize}
        \item Completa la instrucción actual
        \item Guarda el contexto (PC en EPC, copia Status en registro temporal)
        \item Deshabilita interrupciones (bit IE = 0)
    \end{itemize}
    \item \textbf{Identificación de la fuente}: Consulta el registro Cause para determinar qué dispositivo generó la interrupción.
    \item \textbf{Salto a la rutina de servicio}: Carga la dirección base del vector de interrupciones (0x80000180 en MIPS) en PC.
    \item \textbf{Ejecución de ISR}: 
    \begin{itemize}
        \item Guarda registros en pila
        \item Atiende el dispositivo
        \item Limpia la condición de interrupción
    \end{itemize}
    \item \textbf{Recuperación del contexto}: Restaura registros desde la pila.
    \item \textbf{Retorno}: Ejecuta instrucción ERET que:
    \begin{itemize}
        \item Restaura el registro Status
        \item Salta a la dirección guardada en EPC
        \item Reactiva interrupciones
    \end{itemize}
\end{enumerate}

\begin{center}
\begin{tikzpicture}[
    node distance=1.2cm,
    block/.style={rectangle, draw, text width=10em, text centered, rounded corners, minimum height=3em},
    line/.style={draw, -latex', thick}
]
    \node [block] (evento) {Evento en dispositivo};
    \node [block, below=of evento] (solicitud) {Solicitud de interrupción (IRQ)};
    \node [block, below=of solicitud] (deteccion) {Detección por CPU};
    \node [block, below=of deteccion] (contexto) {Guardar contexto (EPC, Status)};
    \node [block, below=of contexto] (identificacion) {Identificar fuente (Cause)};
    \node [block, below=of identificacion] (isr) {Ejecutar ISR};
    \node [block, below=of isr] (retorno) {Restaurar contexto};
    \node [block, below=of retorno] (continuar) {Continuar ejecución (ERET)};

    \path [line] (evento) -- (solicitud);
    \path [line] (solicitud) -- (deteccion);
    \path [line] (deteccion) -- (contexto);
    \path [line] (contexto) -- (identificacion);
    \path [line] (identificacion) -- (isr);
    \path [line] (isr) -- (retorno);
    \path [line] (retorno) -- (continuar);
    
    \draw [line, dashed] ($(continuar.south)+(0,-0.5)$) -- ++(7,0) |- ($(evento.north)+(0,0.5)$) -- (evento.north);
    \node at ($(continuar.south)+(3.5,-0.5)$) [above] {Ciclo continuo};
\end{tikzpicture}
\end{center}

\section*{2. Diferencias entre Gestión de E/S con Sondeo e Interrupciones}

\begin{table}[h]
\centering
\begin{tabularx}{\textwidth}{|>{\raggedright\arraybackslash}X|>{\raggedright\arraybackslash}X|}
\hline
\textbf{Sondeo (Polling)} & \textbf{Interrupciones} \\
\hline
\hline
\begin{minipage}[t]{0.45\textwidth}
\vspace{-\baselineskip}
\begin{itemize}[leftmargin=*,noitemsep]
    \item \textbf{Enfoque}: La CPU verifica periódicamente el estado de los dispositivos
    \item \textbf{Mecanismo}: Ciclo constante de consulta
    \item \textbf{Implementación}:
    \begin{lstlisting}
loop:
    # Verificar estado dispositivo
    lw $t0, DISP_STATUS
    andi $t0, $t0, 1
    beqz $t0, loop
    # Procesar cuando listo
    \end{lstlisting}
\end{itemize}
\end{minipage} & 
\begin{minipage}[t]{0.45\textwidth}
\vspace{-\baselineskip}
\begin{itemize}[leftmargin=*,noitemsep]
    \item \textbf{Enfoque}: Dispositivo notifica a la CPU cuando necesita atención
    \item \textbf{Mecanismo}: Evento asíncrono
    \item \textbf{Implementación}:
    \begin{lstlisting}
main:
    # Trabajo principal
    j main

ISR:
    # Atender interrupción
    eret
    \end{lstlisting}
\end{itemize}
\end{minipage} \\
\hline
\textbf{Ventajas}: & \textbf{Ventajas}: \\
\begin{itemize}[leftmargin=*,noitemsep]
    \item Implementación simple
    \item Comportamiento predecible
    \item Fácil depuración
\end{itemize} &
\begin{itemize}[leftmargin=*,noitemsep]
    \item CPU libre para otras tareas
    \item Bajo consumo energético
    \item Respuesta inmediata a eventos
    \item Manejo eficiente de múltiples dispositivos
\end{itemize} \\
\hline
\textbf{Desventajas}: & \textbf{Desventajas}: \\
\begin{itemize}[leftmargin=*,noitemsep]
    \item Consumo innecesario de ciclos de CPU
    \item Baja eficiencia energética
    \item Latencia variable (depende de frecuencia de sondeo)
    \item No escalable con múltiples dispositivos
\end{itemize} &
\begin{itemize}[leftmargin=*,noitemsep]
    \item Complejidad en implementación
    \item Sobrecarga por cambio de contexto
    \item Posibles problemas de prioridad/race conditions
    \item Dificultad en depuración
\end{itemize} \\
\hline
\end{tabularx}
\end{table}

\subsection*{Comparación detallada}

\subsubsection*{Mecanismo de operación}
\begin{itemize}
    \item \textbf{Sondeo}: La CPU \textbf{inicia} la comunicación verificando activamente el estado del dispositivo en intervalos regulares. Este enfoque es \textbf{orientado a la CPU} y consume recursos constantemente.
    
    \item \textbf{Interrupciones}: El dispositivo \textbf{inicia} la comunicación notificando a la CPU cuando necesita atención. Este enfoque es \textbf{orientado al dispositivo} y permite a la CPU trabajar en otras tareas hasta la notificación.
\end{itemize}

\subsubsection*{Eficiencia de CPU}
\begin{itemize}
    \item \textbf{Sondeo}: Utiliza ciclos de CPU incluso cuando no hay operaciones pendientes. En el peor caso (dispositivo lento), la CPU puede estar inactiva >90\% del tiempo en espera activa.
    
    \item \textbf{Interrupciones}: Permite máximo aprovechamiento de la CPU. Solo consume ciclos cuando realmente hay que atender un dispositivo. En aplicaciones típicas, reduce el uso de CPU en E/S en un 70-90\%.
\end{itemize}

\subsubsection*{Tiempo de respuesta}
\begin{itemize}
    \item \textbf{Sondeo}: Latencia determinada por el intervalo de sondeo. Para garantizar respuesta rápida se requiere alta frecuencia de sondeo, lo que aumenta el consumo de recursos.
    
    \item \textbf{Interrupciones}: Latencia mínima constante. El tiempo entre el evento y el inicio de la ISR es fijo y típicamente pequeño (10-100 ciclos en MIPS).
\end{itemize}

\subsubsection*{Implementación en sistemas MIPS}
\begin{itemize}
    \item \textbf{Sondeo}: No requiere configuración especial de registros de control. Solo necesita acceso directo a registros de estado de dispositivos.
    
    \item \textbf{Interrupciones}: Requiere configuración de registros especiales:
    \begin{itemize}
        \item Habilitación global (bit 0 de Status)
        \item Habilitación por dispositivo (bits 8-15 de Status)
        \item Manejo de EPC y Cause
        \item Implementación de rutinas de servicio (.ktext)
    \end{itemize}
\end{itemize}

\subsubsection*{Casos de uso típicos}
\begin{itemize}
    \item \textbf{Sondeo}: Adecuado para sistemas simples con un solo dispositivo, o cuando la frecuencia de eventos es muy alta (mayor que la capacidad de manejo de interrupciones).
    
    \item \textbf{Interrupciones}: Esencial para sistemas complejos con múltiples dispositivos, sistemas en tiempo real, y aplicaciones donde la eficiencia energética es crítica.
\end{itemize}

\subsubsection*{Ejemplo práctico en buffer.asm}
El programa \texttt{buffer.asm} utiliza interrupciones para:
\begin{itemize}
    \item Atender entrada de teclado solo cuando ocurre una pulsación (no verifica constantemente)
    \item Activar el temporizador para los intervalos de 20 segundos
    \item Combinar ambas interrupciones sin conflicto
\end{itemize}
Implementar esto con sondeo requeriría:
\begin{lstlisting}
main_loop:
    # Verificar temporizador
    lw $t0, TIMER_STATUS
    andi $t0, $t0, 1
    bnez $t0, handle_timer
    
    # Verificar teclado
    lw $t0, KEYBOARD_STATUS
    andi $t0, $t0, 1
    bnez $t0, handle_key
    
    j main_loop
\end{lstlisting}
Este enfoque consumiría 100\% del CPU aunque no hubiera actividad.

\subsubsection*{Ejemplo en semáforo.asm}
El semáforo utiliza interrupciones para:
\begin{itemize}
    \item Detectar pulsaciones del botón (tecla 's') sin verificación constante
    \item Manejar múltiples temporizadores para transiciones de estado
    \item Combinar eventos de teclado y temporizador eficientemente
\end{itemize}

\subsection*{Conclusiones}
La gestión por interrupciones es superior en la mayoría de escenarios modernos debido a su eficiencia y escalabilidad. El sondeo solo se justifica en sistemas muy simples o cuando la frecuencia de eventos es extremadamente alta. En los programas desarrollados (buffer.asm y semáforo.asm), el uso de interrupciones permite una gestión eficiente de múltiples dispositivos y eventos simultáneos con mínimo consumo de recursos.

\section*{3. Ventajas del Uso de Interrupciones en Términos de Uso del Procesador}
\begin{enumerate}
    \item \textbf{Maximización de trabajo útil}: La CPU no desperdicia ciclos verificando dispositivos inactivos
    \item \textbf{Paralelismo efectivo}: Permite solapamiento de cómputo y operaciones de E/S
    \item \textbf{Eficiencia energética}: Permite estados de bajo consumo (sleep) durante esperas
    \item \textbf{Multitarea}: Soporte para múltiples procesos con tiempos de respuesta garantizados
    \item \textbf{Latencia controlada}: Respuesta inmediata a eventos críticos sin retrasos por sondeo
    \item \textbf{Throughput mejorado}: Mayor cantidad de trabajo completado por unidad de tiempo
\end{enumerate}

\section*{4. Registros Especiales para Gestión de Interrupciones en MIPS32}
\begin{table}[h]
\centering
\begin{tabular}{|l|l|p{8cm}|}
\hline
\textbf{Registro} & \textbf{CP0} & \textbf{Función} \\
\hline
Status & Reg 12 & 
\begin{itemize}[leftmargin=*]
    \item Bit 0 (IE): Habilitación global de interrupciones
    \item Bits 8-15 (IM): Máscara para habilitar dispositivos específicos
    \item Bit 1 (EXL): Nivel de excepción (1 = en excepción)
\end{itemize} \\
\hline
Cause & Reg 13 & 
\begin{itemize}[leftmargin=*]
    \item Bits 8-15 (IP): Interrupciones pendientes
    \item Bits 2-6 (ExcCode): Código de excepción
\end{itemize} \\
\hline
EPC & Reg 14 & Almacena la dirección de retorno después de atender interrupción \\
\hline
BadVAddr & Reg 8 & Dirección inválida en fallos de memoria \\
\hline
Context & Reg 4 & Información para manejo de fallos de página \\
\hline
\end{tabular}
\end{table}

\section*{5. Necesidad de Guardar el Contexto en Rutinas de Servicio}
El guardado de contexto es esencial porque:
\begin{itemize}
    \item \textbf{Transparencia}: El programa interrumpido debe continuar como si nada hubiera ocurrido
    \item \textbf{Consistencia de estado}: Registros contienen datos críticos para la ejecución correcta
    \item \textbf{Prevención de corrupción}: La ISR usa registros que podrían estar en uso por el programa principal
    \item \textbf{Multi-anidamiento}: Soporte para interrupciones durante el manejo de otras interrupciones
\end{itemize}

\textbf{Consecuencias de no guardar contexto}:
\begin{itemize}
    \item Corrupción de datos en registros
    \item Comportamiento impredecible del programa
    \item Fallos en cascada
    \item Bloqueos del sistema
\end{itemize}

\section*{6. Excepciones en Sistemas MIPS32}
\subsection*{a) Situaciones generadoras de excepciones}
\begin{enumerate}
    \item Desbordamiento aritmético (overflow)
    \item Instrucción ilegal (código de operación no reconocido)
    \item Fallo de página (dirección virtual sin mapeo físico)
    \item Violación de privilegios (instrucción privilegiada en modo usuario)
    \item Acceso a memoria no alineado (LW/SW en dirección no múltiplo de 4)
    \item Llamada al sistema (SYSCALL)
    \item Breakpoint (BREAK)
    \item Fallo de bus (acceso a dirección física inexistente)
\end{enumerate}

\subsection*{b) Etapas del pipeline que provocan excepciones}
\begin{table}[h]
\centering
\begin{tabular}{|l|p{10cm}|}
\hline
\textbf{Etapa} & \textbf{Excepciones} \\
\hline
IF (Instruction Fetch) & 
\begin{itemize}
    \item Fallo de página en acceso a instrucción
    \item Dirección de instrucción inválida
\end{itemize} \\
\hline
ID (Instruction Decode) & 
\begin{itemize}
    \item Instrucción ilegal
    \item Violación de privilegios
    \item Instrucción no implementada
\end{itemize} \\
\hline
EX (Execute) & 
\begin{itemize}
    \item Desbordamiento aritmético
    \item División por cero
\end{itemize} \\
\hline
MEM (Memory Access) & 
\begin{itemize}
    \item Fallo de página en acceso a datos
    \item Dirección de memoria no alineada
    \item Fallo de bus
\end{itemize} \\
\hline
\end{tabular}
\end{table}

\section*{7. Estrategias para Tratamiento de Excepciones e Interrupciones}
\subsection*{a) Diferencias entre interrupciones y excepciones}
\begin{table}[h]
\centering
\begin{tabular}{|l|l|l|}
\hline
\textbf{Característica} & \textbf{Interrupciones} & \textbf{Excepciones} \\
\hline
Origen & Eventos externos (dispositivos) & Eventos internos (ejecución) \\
\hline
Sincronía & Asíncronas (independientes de instrucción) & Síncronas (relacionadas con instrucción) \\
\hline
Predictibilidad & Ocurren en cualquier momento & Ocurren durante ejecución específica \\
\hline
Ejemplos & Teclado, temporizador & Overflow, página inválida \\
\hline
\end{tabular}
\end{table}

\subsection*{b) Estrategias de tratamiento}
\begin{enumerate}
    \item \textbf{Vectorizado}:
    \begin{itemize}
        \item Diferentes excepciones saltan a distintas direcciones
        \item Implementado mediante tabla de vectores
        \item Más rápido para manejo específico
    \end{itemize}
    
    \item \textbf{No vectorizado}:
    \begin{itemize}
        \item Todas las excepciones saltan a misma dirección (0x80000180)
        \item Software determina causa consultando registro Cause
        \item Implementado en MIPS estándar
    \end{itemize}
\end{enumerate}

\textbf{Función del registro EPC}:
\begin{itemize}
    \item Almacena dirección de la instrucción que causó la excepción
    \item Permite reanudar ejecución después de manejar excepción
    \item Instrucción ERET restaura PC desde EPC
\end{itemize}

\section*{8. Habilitación de Interrupciones}
\subsection*{a) Procedimiento de habilitación}
\textbf{Teclado}:
\begin{lstlisting}
li $v0, 35       # Habilitar interrupciones de teclado
syscall
\end{lstlisting}

\textbf{Pantalla}:
\begin{lstlisting}
li $v0, 36       # Habilitar interrupciones de pantalla
syscall
\end{lstlisting}

\textbf{Procesador} (registro Status):
\begin{lstlisting}
mfc0 $t0, $12    # Leer registro Status
ori $t0, 0x1     # Habilitar globalmente (bit 0)
ori $t0, 0x800   # Habilitar teclado (bit 11)
ori $t0, 0x1000  # Habilitar temporizador (bit 12)
mtc0 $t0, $12    # Escribir registro Status
\end{lstlisting}

\subsection*{b) Consecuencias de habilitar dispositivos pero no el procesador}
\begin{itemize}
    \item Las IRQ se registran en registro Cause (bits pendientes)
    \item CPU ignora completamente las solicitudes
    \item No se ejecuta ninguna rutina de servicio
    \item Dispositivos pueden quedar en estado bloqueado
    \item Eventos críticos no se atienden
    \item Posible pérdida de datos (buffers sobreescritos)
\end{itemize}

\section*{9. Procesamiento de Interrupciones de Reloj}
\subsection*{a) Paso a paso}
\begin{enumerate}
    \item \textbf{Evento}: Temporizador alcanza valor programado
    \item \textbf{Solicitud}: Módulo de temporización activa línea IRQ
    \item \textbf{Detección}: CPU detecta IRQ al final de ciclo de instrucción
    \item \textbf{Guardado hardware}: 
    \begin{itemize}
        \item PC → EPC
        \item Status → registro temporal
        \item IE = 0 (deshabilita interrupciones)
        \item EXL = 1 (entra en nivel excepción)
    \end{itemize}
    \item \textbf{Salto}: PC ← 0x80000180
    \item \textbf{Guardado software}: Rutina salva registros en pila
    \item \textbf{Identificación}: Lee registro Cause para confirmar IRQ de reloj
    \item \textbf{Procesamiento}: 
    \begin{itemize}
        \item Actualizar contadores
        \item Planificar cambio de contexto
        \item Reprogramar temporizador
    \end{itemize}
    \item \textbf{Limpieza}: Resetear flag de interrupción
    \item \textbf{Restauración}: Recupera registros desde pila
    \item \textbf{Retorno}: Ejecuta ERET → restaura Status y PC, reactiva interrupciones
\end{enumerate}

\subsection*{b) Importancia de guardar el contexto}
\begin{itemize}
    \item Garantiza que programa interrumpido continúe correctamente
    \item Previene que ISR corrompa estado de aplicación
    \item Permite múltiples interrupciones anidadas
    \item Soporta rutinas compartidas por diferentes interrupciones
\end{itemize}

\section*{10. Interrupciones de Reloj para Control de Ejecución}
\subsection*{a) Aplicaciones}
\begin{enumerate}
    \item \textbf{Prevenir bucles infinitos}:
    \begin{itemize}
        \item Programar temporizador con límite de tiempo
        \item En ISR: terminar proceso o tomar acción correctiva
        \begin{lstlisting}
# Configurar temporizador
li $a0, MAX_TIME
li $v0, 33
syscall
        \end{lstlisting}
    \end{itemize}
    
    \item \textbf{Finalizar programas que superan tiempo máximo}:
    \begin{itemize}
        \item Esencial en sistemas de tiempo real
        \item ISR fuerza finalización si se excede tiempo
    \end{itemize}
\end{enumerate}

\subsection*{b) Manejo de programa que finaliza antes}
\begin{itemize}
    \item Deshabilitar interrupción: Cancelar temporizador pendiente
    \item Reiniciar temporizador: Para siguiente intervalo
    \item Liberar recursos: Memoria, dispositivos asociados
    \item Notificar sistema: Registrar finalización exitosa
\end{itemize}

\section*{11. Análisis y Discusión de Resultados}
\subsection*{a) Buffer de Entrada de Teclado (buffer.asm)}
\textbf{Implementación}:
\begin{itemize}
    \item Buffer circular de 100 caracteres
    \item Head y tail para gestión FIFO
    \item Interrupciones de teclado y temporizador
\end{itemize}

\textbf{Logros}:
\begin{itemize}
    \item Manejo eficiente de entrada asíncrona
    \item Sin pérdida de datos durante impresión
    \item Uso óptimo de CPU durante esperas
\end{itemize}

\textbf{Mejoras potenciales}:
\begin{itemize}
    \item Implementar lock para acceso concurrente
    \item Añadir manejo de buffer lleno
    \item Soporte para teclas especiales (Backspace)
\end{itemize}

\subsection*{b) Semáforo con Pulsador (semaforo.asm)}
\textbf{Implementación}:
\begin{itemize}
    \item Máquina de estados: Verde → Amarillo → Rojo
    \item Control por pulsador (tecla 's')
    \item Temporizadores para transiciones
\end{itemize}

\textbf{Logros}:
\begin{itemize}
    \item Respuesta inmediata a evento externo
    \item Transiciones temporizadas precisas
    \item Simulación completa de comportamiento real
\end{itemize}

\textbf{Limitaciones}:
\begin{itemize}
    \item No manejo de pulsaciones múltiples
    \item Sin prioridad para eventos concurrentes
    \item Estados bloqueantes durante temporizaciones
\end{itemize}

\subsection*{c) Conclusiones Generales}
\begin{itemize}
    \item Las interrupciones permiten construir sistemas reactivos eficientes
    \item MIPS provee mecanismos robustos para gestión de eventos
    \item El diseño modular facilita expansión y mantenimiento
    \item La gestión adecuada de contexto es crítica para estabilidad
    \item Los temporizadores son esenciales para sistemas en tiempo real
    \item La combinación de interrupciones de hardware y temporizadores permite implementar sistemas complejos con comportamiento determinista
\end{itemize}

\end{document}